\documentclass[12pt]{article}
\usepackage[utf8]{inputenc}
\usepackage[T1]{fontenc}
\usepackage{hyperref}
\title{Tjurruru Language Resources for K-12 Curriculum Development}
\author{Compiled by Claude Sonnet 5 (claude-sonnet-5) for Andr\'as Kornai}
\date{}
\begin{document}
\maketitle
\section{Introduction}
Tjurruru (also spelled Jurruru or Tjuroro; ISO 639-3 tju) is an extinct Aboriginal language of the Pilbara region of Western Australia. The Tjuroro people lived along the Ashburton River. One source says two speakers were left in 1967 and none by 1986. A search for a Queensland wordlist found nothing for this language. The documentation is very thin.
\section{Primary Educational Resources}
None found.
\section{Linguistic and Reference Materials}
\begin{itemize}
\item AIATSIS AUSTLANG record (code W33). \url{https://collection.aiatsis.gov.au/austlang/language/w33}
\item Wikipedia, ``Jurruru language''. \url{https://wikipedia.com/wiki/ISO_639:tju}
\item Wikipedia, ``Tjuroro people''. \url{https://en.wikipedia.org/wiki/Tjuroro_people}
\item Cosmos entry on Jurruru (PDF). \url{https://www.hgcosmos.org/assets/culture/AZ/J/Jurruru.pdf}
\end{itemize}
The search summary also named Thieberger (1993), \emph{Handbook of Western Australian Aboriginal Languages South of the Kimberley Region}, as the main source of word-level documentation. I did not see a URL for it.
\section{Religious and Community Resources}
None found.
\section{Digital and Audio Resources}
None found.
\section{Curriculum Development Context}
A heritage curriculum would have to rest on the small amount of archival material and on the descendant community's own wishes. The language is not in use, so any revitalization would be reconstruction from limited records.
\end{document}